\SourceLeaf{357}

\noindent when we substitute $X = m'm''$, $X' = mm''$, $X'' = mm'$, while six times
the volume of the pyramid thus formed with the origin is $= mm'm''\BareRoot{D}$.
Hence the perpendicular [is]
\[
= \BareRoot{\dfrac{D}{F\left(\dfrac{1}{m}, \dfrac{1}{m'}, \dfrac{1}{m''}\right)}};
\]

$T$, $T'$, $T''$ stand in the same relation to the form $\begin{pmatrix} AD, A'D, A''D \\ BD, B'D, B''D \end{pmatrix}$
as $t$, $t'$, $t''$ do to $\begin{pmatrix} a, a', a'' \\ b, b', b'' \end{pmatrix}$.



\Needspace{12\baselineskip}
\begin{center}
\bfseries REMARK.
\end{center}

The preceding note is in a notebook and dates from July 1831. It
shows that \textsc{Gauss} also investigated complex quantities with more than two units.
In the concluding remark of his own notice of Theoria residuorum
biquadraticorum, Commentatio secunda, on 23 April 1831 (vol. II, pp. 169 to 198), he had already mentioned such quantities.
This is why the note is included again here, although it is already printed in vol. II, pp. 305 to
307 because of its relation to the theory of ternary forms.
\begin{flushright}
\textsc{Stäckel.}
\end{flushright}


